% !TEX program = xelatex
% ============================================================
%  مواصفات تقنية — Technical Specification Template
%  قالب لوثيقة المواصفات التقنية للمشاريع الهندسية والبرمجية.
%  Compile with: xelatex main.tex (run twice)
% ============================================================
%  Features:
%    - Version and status table at top
%    - Overview, functional/non-functional requirements tables
%    - Architecture block diagram (TikZ placeholder)
%    - Data model, API interfaces table
%    - Constraints, assumptions, open questions
% ============================================================

\documentclass[11pt,a4paper]{article}

\usepackage[a4paper,margin=2.5cm,top=2.5cm,bottom=2.5cm,headheight=16pt]{geometry}
\usepackage{graphicx}
\usepackage{array}
\usepackage{booktabs}
\usepackage{tabularx}
\usepackage{multirow}
\usepackage{enumitem}
\usepackage{float}
\usepackage{caption}
\usepackage{titlesec}
\usepackage{fancyhdr}
\usepackage{setspace}
\usepackage{xcolor}
\usepackage{amsmath}
\usepackage{amssymb}
\usepackage{tikz}
\usetikzlibrary{positioning,calc,shapes.geometric,arrows.meta}

\usepackage{bidi}
\usepackage[no-math]{fontspec}
\usepackage[quiet]{polyglossia}

\setmainfont[Scale=1]{NotoKufiArabic-Regular.ttf}
\newfontfamily\arabicfont[Script=Arabic,Scale=0.95]{NotoKufiArabic-Regular.ttf}
\newfontfamily\englishfont[Scale=0.95]{TeX Gyre Termes}

\setdefaultlanguage[calendar=gregorian,hijricorrection=1]{arabic}
\setotherlanguage[variant=british]{english}

\usepackage[breaklinks,hidelinks]{hyperref}
\hypersetup{pdftitle={مواصفات تقنية},pdfauthor={اسم الشركة}}

\addto\captionsarabic{%
  \renewcommand{\contentsname}{الفهرس}%
  \renewcommand{\figurename}{الشكل}%
  \renewcommand{\tablename}{الجدول}%
}

\titleformat{\section}{\Large\bfseries}{\thesection}{0.7em}{}
\titleformat{\subsection}{\large\bfseries}{\thesubsection}{0.6em}{}
\titlespacing*{\section}{0pt}{1.4ex plus 0.4ex}{0.9ex plus 0.3ex}

\pagestyle{fancy}
\fancyhf{}
\fancyhead[R]{\small\itshape مواصفات تقنية (Technical Spec)}
\fancyhead[L]{\small اسم المنتج}
\fancyfoot[C]{\small\thepage}
\renewcommand{\headrulewidth}{0.3pt}

\setlist[itemize]{topsep=0.3em, itemsep=0.15em, parsep=0pt}
\setlist[enumerate]{topsep=0.3em, itemsep=0.15em, parsep=0pt}

\newcolumntype{L}[1]{>{\raggedright\arraybackslash}p{#1}}
\newcolumntype{C}[1]{>{\centering\arraybackslash}p{#1}}
\newcolumntype{R}[1]{>{\raggedleft\arraybackslash}p{#1}}

\definecolor{headerColor}{HTML}{1A3C6B}
\definecolor{lightRow}{HTML}{F0F4F8}
\definecolor{blockFill}{HTML}{E8EEF4}
\definecolor{accentColor}{HTML}{2E6DA4}

\setstretch{1.3}

\begin{document}

% ============================================================
% TITLE BLOCK
% ============================================================
\begin{center}
{\LARGE\bfseries مواصفات تقنية (Technical Specification)}\\[0.5em]
{\Large اسم المنتج --- وثيقة المواصفات}\\[1em]
\end{center}

% ---- Version and status table ----
\begin{table}[H]
\centering
\renewcommand{\arraystretch}{1.4}
\begin{tabularx}{\textwidth}{L{3.5cm} X L{3.5cm} L{3cm}}
\toprule
\textbf{اسم المنتج:} & \multicolumn{3}{l}{---} \\
\textbf{الإصدار:} & --- & \textbf{الحالة:} & مسودة (Draft) \\
\textbf{اسم الشركة:} & --- & \textbf{التاريخ:} & --- \\
\textbf{المؤلف:} & --- & \textbf{المراجع:} & --- \\
\bottomrule
\end{tabularx}
\end{table}

\vspace{0.5em}

% ============================================================
% 1. OVERVIEW
% ============================================================
\section{نظرة عامة (Overview)}
\label{sec:overview}

% TODO: اكتب نظرة عامة مختصرة على المنتج والغرض من هذه المواصفات.
هذا نص تجريبي للنظرة العامة. يقدم وصفاً مختصراً للمنتج والغرض من وثيقة المواصفات
التقنية، بالإضافة إلى الجمهور المستهدف ونطاق التطبيق.

\subsection{الغرض (Purpose)}
% TODO: اذكر الغرض من الوثيقة.
وصف الغرض من وثيقة المواصفات والقرارات التي تدعمها.

\subsection{النطاق (Scope)}
% TODO: اذكر نطاق المنتج وما يغطيه.
يحدد هذا القسم نطاق المنتج والمكونات المشمولة وغير المشمولة.

% ============================================================
% 2. REQUIREMENTS
% ============================================================
\section{المتطلبات (Requirements)}
\label{sec:requirements}

\subsection{المتطلبات الوظيفية (Functional Requirements)}
% TODO: املأ جدول المتطلبات الوظيفية.
\begin{table}[H]
\centering
\caption{المتطلبات الوظيفية}
\renewcommand{\arraystretch}{1.3}
\begin{tabularx}{\textwidth}{C{1.2cm} L{4cm} X L{2.5cm}}
\toprule
\textbf{المعرف} & \textbf{المتطلب} & \textbf{الوصف} & \textbf{الأولوية} \\
\midrule
FR-01 & متطلب وظيفي أول & وصف مختصر للمتطلب الوظيفي الأول. & عالية \\
FR-02 & متطلب وظيفي ثانٍ & وصف مختصر للمتطلب الوظيفي الثاني. & عالية \\
FR-03 & متطلب وظيفي ثالث & وصف مختصر للمتطلب الوظيفي الثالث. & متوسطة \\
FR-04 & متطلب وظيفي رابع & وصف مختصر للمتطلب الوظيفي الرابع. & منخفضة \\
\bottomrule
\end{tabularx}
\end{table}

\subsection{المتطلبات غير الوظيفية (Non-Functional Requirements)}
% TODO: املأ جدول المتطلبات غير الوظيفية.
\begin{table}[H]
\centering
\caption{المتطلبات غير الوظيفية}
\renewcommand{\arraystretch}{1.3}
\begin{tabularx}{\textwidth}{C{1.2cm} L{3cm} L{2.5cm} X}
\toprule
\textbf{المعرف} & \textbf{الفئة} & \textbf{المقياس} & \textbf{الوصف} \\
\midrule
NFR-01 & الأداء (Performance) & زمن الاستجابة & زمن استجابة أقل من \LR{200 ms}. \\
NFR-02 & الموثوقية (Reliability) & التوفر & توفر لا يقل عن \LR{99.9\%}. \\
NFR-03 & الأمان (Security) & التشفير & تشفير البيانات أثناء النقل (\LR{TLS 1.2+}). \\
NFR-04 & قابلية التوسع (Scalability) & الحمل & دعم \LR{10{,}000} مستخدم متزامن. \\
\bottomrule
\end{tabularx}
\end{table}

% ============================================================
% 3. ARCHITECTURE
% ============================================================
\section{المعمارية (Architecture)}
\label{sec:architecture}

% TODO: استبدل مخطط المعمارية المؤقت بمخطط المشروع الفعلي.
يوضح الشكل التالي المعمارية العامة للنظام على شكل كتل وظيفية.

\begin{figure}[H]
\centering
\begin{tikzpicture}[
  every node/.style={font=\small},
  block/.style={fill=blockFill,draw=accentColor,thick,rounded corners=3pt,inner sep=6pt,minimum width=2.6cm,minimum height=1.2cm,align=center},
  arrow/.style={draw=accentColor,thick,-{Stealth[length=4mm]}},
]
\node[block] (client)  at (0,0)    {العميل\\(\LR{Client})};
\node[block] (api)     at (4,0)    {بوابة \LR{API}\\(\LR{API Gateway})};
\node[block] (service) at (8,0)    {الخدمات\\(\LR{Services})};
\node[block] (db)      at (8,-2.5) {قاعدة البيانات\\(\LR{Database})};
\node[block] (cache)   at (4,-2.5) {الذاكرة المؤقتة\\(\LR{Cache})};

\draw[arrow] (client)  -- (api);
\draw[arrow] (api)     -- (service);
\draw[arrow] (service) -- (db);
\draw[arrow] (api)     -- (cache);
\end{tikzpicture}
\caption{مخطط المعمارية العامة (Architecture Block Diagram)}
\end{figure}

% ============================================================
% 4. DATA MODEL
% ============================================================
\section{نموذج البيانات (Data Model)}
\label{sec:data-model}

% TODO: صف الكيانات الرئيسية وعلاقاتها.
\begin{table}[H]
\centering
\caption{الكيانات الرئيسية ونموذج البيانات}
\renewcommand{\arraystretch}{1.3}
\begin{tabularx}{\textwidth}{L{3cm} L{3cm} X}
\toprule
\textbf{الكيان (Entity)} & \textbf{المعرف (ID)} & \textbf{الحقول الرئيسية (Key Fields)} \\
\midrule
المستخدم (\LR{User}) & \LR{user\_id} & \LR{name, email, role, created\_at} \\
المنتج (\LR{Product}) & \LR{product\_id} & \LR{name, price, stock, category} \\
الطلب (\LR{Order}) & \LR{order\_id} & \LR{user\_id, total, status, items} \\
\bottomrule
\end{tabularx}
\end{table}

% ============================================================
% 5. INTERFACES (API)
% ============================================================
\section{الواجهات (Interfaces / API)}
\label{sec:interfaces}

% TODO: املأ جدول نقاط نهاية الـ API.
\begin{table}[H]
\centering
\caption{نقاط نهاية واجهة البرمجة (API Endpoints)}
\renewcommand{\arraystretch}{1.3}
\begin{tabularx}{\textwidth}{L{4.5cm} C{2cm} X}
\toprule
\textbf{المسار (Endpoint)} & \textbf{الطريقة (Method)} & \textbf{الوصف (Description)} \\
\midrule
\LR{/api/v1/users}      & \LR{GET}  & استرجاع قائمة المستخدمين. \\
\LR{/api/v1/users}      & \LR{POST} & إنشاء مستخدم جديد. \\
\LR{/api/v1/users/\{id\}} & \LR{GET}  & استرجاع تفاصيل مستخدم. \\
\LR{/api/v1/orders}     & \LR{GET}  & استرجاع قائمة الطلبات. \\
\bottomrule
\end{tabularx}
\end{table}

% ============================================================
% 6. CONSTRAINTS
% ============================================================
\section{القيود (Constraints)}
\label{sec:constraints}

% TODO: اذكر القيود التقنية والتشغيلية.
\begin{itemize}
    \item قيد تقني: التوافق مع \LR{Python 3.10+} و \LR{Node.js 18+}.
    \item قيد تشغيلي: التشغيل في بيئة \LR{Linux} فقط.
    \item قيد تنظيمي: الالتزام بمعايير \LR{GDPR} لحماية البيانات.
    \item قيد الموارد: ميزانية استضافة لا تتجاوز \LR{\$5{,}000} سنوياً.
\end{itemize}

% ============================================================
% 7. ASSUMPTIONS
% ============================================================
\section{الافتراضات (Assumptions)}
\label{sec:assumptions}

% TODO: اذكر الافتراضات التي بنيت عليها المواصفات.
\begin{itemize}
    \item افتراض أن المستخدمين لديهم اتصال إنترنت مستقر.
    \item افتراض أن فريق التطوير يضم \LR{5} مهندسين على الأقل.
    \item افتراض توفر بنية تحتية سحابية (\LR{Cloud}) جاهزة.
\end{itemize}

% ============================================================
% 8. OPEN QUESTIONS
% ============================================================
\section{الأسئلة المفتوحة (Open Questions)}
\label{sec:open-questions}

% TODO: سجل الأسئلة التي لم تُحسم بعد.
\begin{table}[H]
\centering
\caption{الأسئلة المفتوحة}
\renewcommand{\arraystretch}{1.3}
\begin{tabularx}{\textwidth}{C{1cm} X L{3cm} C{2cm}}
\toprule
\textbf{رقم} & \textbf{السؤال} & \textbf{المسؤول} & \textbf{الحالة} \\
\midrule
1 & هل سيتم دعم تعدد اللغات في الإصدار الأول؟ & --- & مفتوح \\
2 & ما هي سياسة الاحتفاظ بالبيانات (\LR{Retention})؟ & --- & مفتوح \\
3 & هل هناك تكامل مع أنظمة خارجية؟ & --- & مفتوح \\
\bottomrule
\end{tabularx}
\end{table}

\end{document}
